\documentclass[10pt]{beamer}
\usetheme{Madrid}
\usecolortheme{default}
\usepackage[utf8]{inputenc}
\usepackage[spanish]{babel}
\usepackage{amsmath}
\usepackage{amsfonts}
\usepackage{amssymb}
\usepackage{graphicx}
\usepackage{booktabs}
\usepackage{hyperref}

% --- Metadata ---
\title[MCSC Solver]{MCSC Resolver: TalentBridge Connect}
\subtitle{Búsqueda de Equipos de Costo Mínimo con Cobertura de Habilidades}
\author[DAA 2026]{Diseño y Análisis de Algoritmos}
\institute[DAA]{Facultad de Ingeniería}
\date{\today}

\begin{document}

% 1. Portada
\begin{frame}
    \titlepage
\end{frame}

% 2. Introducción
\begin{frame}{Introducción y Motivación}
    \begin{itemize}
        \item \textbf{Contexto:} TalentBridge Connect es una plataforma que conecta proyectos con freelancers.
        \item \textbf{El Problema:} Seleccionar el equipo más económico que cubra todos los requisitos técnicos del proyecto.
        \item \textbf{Desafío:} Los freelancers tienen múltiples habilidades con distintos niveles y costos variados.
        \item \textbf{Objetivo:} Desarrollar un motor de búsqueda eficiente y escalable para optimizar la selección.
    \end{itemize}
\end{frame}

% 3. Definición del Problema (MCSC)
\begin{frame}{Formalización del Problema (MCSC)}
    Definimos el \textbf{Minimum-Cost Skill-Cover} como:
    \vspace{0.5cm}
    \begin{block}{Parámetros}
        \begin{itemize}
            \item $F = \{f_1, f_2, \dots, f_m\}$ freelancers con costo $c_i$.
            \item $H = \{h_1, h_2, \dots, h_n\}$ habilidades.
            \item Cada $f_i$ posee $h_j$ con nivel $L_{ij}$.
            \item El proyecto requiere habilidades $h_j$ con nivel mínimo $R_j$.
        \end{itemize}
    \end{block}
    \begin{block}{Objetivo}
        Minimizar $\sum_{i=1}^{m} c_i x_i$, donde $x_i \in \{0, 1\}$, asegurando que para cada requisito $j$, exista al menos un $f_i$ con $x_i=1$ tal que $L_{ij} \ge R_j$.
    \end{block}
\end{frame}

% 4. Análisis de Complejidad
\begin{frame}{Análisis de Complejidad}
    \begin{itemize}
        \item \textbf{Clasificación:} NP-Completo.
        \item \textbf{Demostración:} Reducción desde el problema clásica \textbf{Set Cover}.
        \item \textbf{Intuición:} Si todos los costos fueran 1 y niveles mínimos fueran 1, MCSC es idéntico a Set Cover.
        \item \textbf{Implicación:} No esperamos encontrar un algoritmo óptimo polinomial. Necesitamos estrategias inteligentes.
    \end{itemize}
\end{frame}

% 5. Metodología y Diseño
\begin{frame}{Enfoques Implementados}
    Para abordar la complejidad, se diseñaron 5 estrategias:
    \begin{enumerate}
        \item \textbf{Fuerza Bruta:} Backtracking con poda avanzada para óptimo exacto.
        \item \textbf{Greedy:} Algoritmo voraz basado en ratio costo-beneficio.
        \item \textbf{Local Search:} Mejora iterativa con operadores \textbf{2-1 Swap}.
        \item \textbf{Algoritmo Genético:} Evolución poblacional para búsqueda global.
        \item \textbf{Híbrido:} Preprocesado matemático + selección adaptativa de algoritmo.
    \end{enumerate}
\end{frame}

% 6. Heurísticas Clásicas
\begin{frame}{Algoritmo Greedy (Voraz)}
    \begin{itemize}
        \item Selecciona en cada paso al freelancer que maximiza:
        \[ \text{Ratio} = \frac{\text{Nuevas habilidades cubiertas}}{\text{Costo}} \]
        \item \textbf{Ventaja:} Extremadamente rápido (milisegundos).
        \item \textbf{Desventaja:} Se queda atrapado en óptimos locales. En el caso $m=500$, obtuvo un costo de 81.0 frente al óptimo de 67.0 ($20\%$ de error).
    \end{itemize}
\end{frame}

% 7. Búsqueda Local Refinada
\begin{frame}{Búsqueda Local con 2-1 Swap}
    Mejora una solución inicial mediante cambios en el vecindario:
    \begin{itemize}
        \item \textbf{Eliminación Redundante:} Quitar freelancers innecesarios.
        \item \textbf{Intercambio 2-1 (Swap):} Reemplazar dos freelancers por uno solo más barato que cubra sus funciones.
        \item \textbf{Intercambio 1-1:} Sustituir por un perfil más económico.
    \end{itemize}
    \begin{exampleblock}{Resultado}
        Este operador permitió que LS alcanzara el óptimo absoluto (67.0) en el caso $m=500$, superando incluso al genético estándar.
    \end{exampleblock}
\end{frame}

% 8. Metaheurística Evolutiva
\begin{frame}{Algoritmo Genético}
    Simula la evolución para explorar el espacio de soluciones:
    \begin{itemize}
        \item \textbf{Población:} 50 individuos representados por vectores binarios.
        \item \textbf{Selección:} Torneo para preservar los mejores genes.
        \item \textbf{Cruce y Mutación:} Uniformes para diversificar la búsqueda.
        \item \textbf{Reparación:} Mecanismo Greedy para asegurar que cada individuo sea una solución válida.
    \end{itemize}
\end{frame}

% 9. La Clave: Kernelization
\begin{frame}{Kernelization (Reducción de Instancia)}
    Antes de optimizar, simplificamos el problema mediante reglas lógicas:
    \begin{enumerate}
        \item \textbf{Dominancia:} Si $f_A$ es más barato y cubre más que $f_B$, descartamos $f_B$.
        \item \textbf{Esencialidad:} Si una habilidad solo la tiene un freelancer, este debe ser incluido de forma obligatoria.
    \end{enumerate}
    \begin{alertblock}{Impacto}
        En la instancia $m=150$, las reglas redujeron el problema a solo **44 candidatos**, permitiendo el uso de Fuerza Bruta en un entorno masivo.
    \end{alertblock}
\end{frame}

% 10. Estrategia Híbrida Adaptativa
\begin{frame}{Solucionador Híbrido}
    El sistema decide la mejor herramienta tras el preprocesado:
    \centering
    \begin{itemize}
        \item \textbf{Entrada:} Instancia original $F$.
        \item \textbf{Fase 1:} Aplicar Kernelization $\to F'$.
        \item \textbf{Fase 2 (Decisión):}
        \begin{itemize}
            \item Si $|F'| \le 750$: Fuerza Bruta (Garantía de óptimo).
            \item Si $|F'| > 750$: Algoritmo Genético + Local Search.
        \end{itemize}
    \end{itemize}
    \vspace{0.3cm}
    \textit{"Lo mejor de los dos mundos: exactitud donde es posible, metaheurística donde es necesario."}
\end{frame}

% 11. Resultados: Calidad y Tiempo
\begin{frame}{Análisis Experimental}

            \textbf{Costos (Solución Calidad)}
            \resizebox{\textwidth}{!}{
	\begin{tabular}{@{}lccccc@{}}
	\toprule
	\textbf{m} & \textbf{Brute-Force} & \textbf{Greedy} & \textbf{Local Search} & \textbf{Genetic} & \textbf{Hybrid} \\
	\midrule
	10 & 52.0 & 52.0 & 52.0 & 52.0 & 52.0 \\
	20 & 149.0 & 149.0 & 149.0 & 149.0 & 149.0 \\
	50 & 96.0 & 116.0 & 96.0 & 96.0 & 96.0 \\
	150 & 56.0 & 56.0 & 56.0 & 56.0 & 56.0 \\
	300 & 67.0 & 67.0 & 67.0 & 67.0 & 67.0 \\
	500 & 67.0 & 81.0 & 67.0 & 73.0 & 67.0 \\
	750 & 65.0 & 84.0 & 74.0 & 65.0 & 65.0 \\
	1000 & N/A & 83.0 & 68.0 & 86.0 & 81.0 \\
	\bottomrule
\end{tabular}
            }
    

\end{frame}

\begin{frame}{Análisis Experimental}
	

			\textbf{Tiempos (Segundos)}
			\resizebox{\textwidth}{!}{
	\begin{tabular}{@{}lccccc@{}}
	\toprule
	\textbf{m} & \textbf{Brute-Force} & \textbf{Greedy} & \textbf{Local Search} & \textbf{Genetic} & \textbf{Hybrid} \\
	\midrule
	10 & 0.0001 & 0.0000 & 0.0000 & 0.0166 & 0.0001 \\
	20 & 0.0005 & 0.0000 & 0.0001 & 0.0211 & 0.0002 \\
	50 & 0.0028 & 0.0001 & 0.0002 & 0.0223 & 0.0007 \\
	150 & 0.0506 & 0.0003 & 0.0004 & 0.0300 & 0.0039 \\
	300 & 0.7608 & 0.0011 & 0.0015 & 0.0533 & 0.2966 \\
	500 & 4.6766 & 0.0033 & 0.0040 & 0.0673 & 2.4421 \\
	750 & 18.7862 & 0.0054 & 0.0067 & 0.1189 & 15.4450 \\
	1000 & N/A & 0.0213 & 0.0360 & 0.4829 & 0.8506 \\
	\bottomrule
\end{tabular}
			}

\end{frame}

% 12. Comparativa Crítica
\begin{frame}{Discusión de Resultados}
    \begin{itemize}
        \item \textbf{Superioridad del Híbrido:} Es el único que mantiene el óptimo (65.0 en $m=750$) gracias a que el preprocesado permite usar FB.
        \item \textbf{Poder del Local Search:} Con el operador 2-1 Swap, es el algortimo más eficiente para escalas masivas ($m=1000$), logrando el costo más bajo (68.0).
        \item \textbf{Limitación del Greedy:} Su naturaleza "miope" lo desvía significativamente en problemas densos.
        \item \textbf{Escalabilidad:} El sistema maneja 1000 perfiles en menos de 1 segundo con alta calidad.
    \end{itemize}
\end{frame}

% 13. Conclusiones y Futuro
\begin{frame}{Conclusiones}
    \begin{itemize}
        \item La **Kernelization** es el factor más determinante para habilitar exactitud en problemas NP-Duros.
        \item La estrategia **Híbrida** ofrece la mejor garantía de calidad para una plataforma real.
        \item Para sistemas de alto volumen y tiempo real, el **Local Search refinado** es la mejor opción.
    \end{itemize}
    \textbf{Trabajo Futuro:}
    \begin{itemize}
        \item Incorporar restricciones de confiabilidad y disponibilidad.
        \item Explorar reglas de reducción basadas en Programación Lineal (LP Relaxation).
    \end{itemize}
\end{frame}

\begin{frame}
	\titlepage
\end{frame}

\end{document}
